% RESUMO--------------------------------------------------------------------------------

\begin{resumo}[RESUMO]
\begin{SingleSpacing}

\imprimirautorcitacao. \imprimirtitulo. \imprimirdata. \pageref{LastPage} f. \imprimirprojeto\ -- \imprimirprograma, \imprimirinstituicao. \imprimirlocal, \imprimirdata.\\

O emprego formal em Tecnologia da Informação (TI) no Brasil cresceu de forma acelerada nas últimas duas décadas, mas as análises disponíveis sobre o setor costumam se apoiar em agregados fechados ou em recortes temporais curtos, sem permitir reprodução. Este trabalho construiu uma plataforma analítica completa sobre os microdados da Relação Anual de Informações Sociais (RAIS) e do Cadastro Geral de Empregados e Desempregados (CAGED), cobrindo 2007 a 2026, organizada em arquitetura Medalhão (Bronze, Silver e Gold) executada inteiramente em nuvem: o \textit{data lake} reside em repositórios públicos do Hugging Face Hub, e o processamento é feito por \textit{workflows} do GitHub Actions encadeados em um grafo de dependências, que verifica periodicamente o servidor público do Ministério do Trabalho e Emprego e reprocessa apenas os anos com dado novo. O desenvolvimento seguiu o CRISP-DM, cujo ciclo levou a novas fontes, à revisão do recorte e à reformulação das perguntas. Os arquivos Apache Parquet são consultados pelo motor \textit{in-process} DuckDB sem importação prévia e apresentados em um \textit{dashboard} \textit{web} em React, publicado como site estático, no qual o DuckDB-Wasm executa as consultas no navegador do visitante. Foram ingeridos 661~GB de microdados em texto, guardados em 56,4~GB de Parquet, dos quais o recorte de tecnologia, definido pela união de atividade econômica (CNAE) e ocupação (CBO), resultou em 23,8 milhões de vínculos da RAIS e 12,1 milhões de movimentações do CAGED. A reprodutibilidade foi verificada recontando o recorte no dado bruto em todos os anos da RAIS e do CAGED, o que apontou e corrigiu uma duplicata de 2007, e terminou com diferença zero, e os totais da lente de ocupação ficaram a cerca de 1\% dos de um painel público do Ministério. A auditoria encontrou também o funcionalismo estadual de Rondônia declarado como programador de sistemas em oito anos da série, defeito que, corrigido, mudou a taxa de crescimento do setor e invertia a tendência do hiato de gênero. Os resultados mostram que o estoque de vínculos formais de TI passou de 519.271 em 2007 para 1.361.231 em 2025, crescimento de 162\%, ou 5,5\% ao ano, em ondas, com 38\% do ganho concentrado em 2021 e 2022; que 58\% dos profissionais de TI trabalham fora de empresas de TI, mas eram 71\% em 2007, de modo que o crescimento se concentrou nas próprias empresas de TI; que a remuneração mediana medida em salários mínimos ficou estável (3,39 para 3,41); que o hiato salarial racial caiu, por melhora de acesso, enquanto o de gênero surgiu depois de 2015 e chegou a 7,8\%, inteiramente não explicado por características observáveis; e que um estimador de \textit{nowcast} do estoque, validado retendo cada ano de 2014 a 2025, errou em média 1,1\%, contra 5,1\% do estimador ingênuo. Na previsão do saldo mensal, o SARIMA selecionado por validação de origem móvel superou o \textit{baseline} ingênuo sazonal em 33\% de erro absoluto médio, contra 10,5\% do melhor Prophet. Como o \textit{data lake} é público, as camadas tratadas são também dados abertos reutilizáveis por terceiros.\\

\textbf{Palavras-chave}: mercado de trabalho em TI; engenharia de dados; RAIS; CAGED; dados abertos.

\end{SingleSpacing}
\end{resumo}
